\chapter{Cosmic web, tidal fields and morphology}\label{ch:cosmic-web}
\chapterstatus{written}{2026-11-14}

The large-scale structure is classified into voids, sheets, filaments and knots from the eigenvalues of the tidal and
velocity-shear tensors (T-web, V-web) or from the stream crossings of the particles (ORIGAMI), and its morphology is measured
with Minkowski functionals of excursion sets.

\begin{modulebox}{CosmicWeb, TidalTensor, Morphology}
\nsp{FML::COSMICWEB}: \file{CosmicWeb/CosmicWeb.h} (\code{TWeb}, \code{VWeb}, \code{classify_tensor}, \cls{TWebOptions}),
\file{CosmicWeb/ORIGAMI.h} (\code{ORIGAMI}), \file{CosmicWeb/ORIGAMIHalos.h} (\code{ORIGAMIHalos}); \file{TidalTensor/Hessian.h}
(\code{ComputeHessianWithFT}, \code{SymmetricTensorEigensystem}); \nsp{FML::MORPHOLOGY}: \file{Morphology/VoxelMorphology.h}
(\code{excursions}, \code{measure}).\par
Documentation: \file{CosmicWeb/README.md}, \file{TidalTensor/Workspace.md}, \file{Morphology/README.md}. Tests:
\file{tests/test_cosmicweb.cpp}, \file{tests/test_cosmicweb_pencil.cpp}, \file{tests/test_reconstruction.cpp}, \file{Morphology/tests},
\file{tests/test_morphology_pencil.cpp}.
\end{modulebox}

\section{T-web and V-web}\label{sec:web-tweb}

The T-web \cite{Hahn2007,ForeroRomero2009} uses the tidal tensor of the (Gaussian-smoothed) density contrast,
\begin{equation}
T_{ij}=\partial_i\partial_j\phi,\qquad\nabla^2\phi=\delta,\qquad T_{ij}(\bm k)=\frac{k_ik_j}{k^2}\delta(\bm k),
\end{equation}
whose trace is $\delta$; the V-web \cite{Hoffman2012} the normalised velocity shear
\begin{equation}
\Sigma_{ij}=-\frac{\partial_iv_j+\partial_jv_i}{2H},
\end{equation}
e.g.\ from the DTFE velocity field (\cref{sec:voids-dtfe}). A cell is a void, sheet, filament or knot if 0, 1, 2 or 3 eigenvalues exceed the
threshold $\lambda_{\rm th}$ (positive values, $\sim0.1$--0.4, give visually better webs \cite{ForeroRomero2009,Hoffman2012}). The eigenvalues
of the symmetric $3\times3$ matrices are computed by cyclic Jacobi rotations, accurate also for (nearly) degenerate eigenvalues, without GSL.
The result holds the classification, optionally the eigenvalues, and the volume and mass fractions of each class; the labels can be stored as one
byte per cell (\code{ClassificationStorage::Bytes}: 1\,GiB instead of 8\,GiB at $1024^3$). \code{TWeb} also takes particles and assigns them to a
grid first. Slabs and pencils are supported.

\begin{cppcode}[T-web classification]
#include <FML/CosmicWeb/CosmicWeb.h>
using namespace FML::COSMICWEB;
TWebOptions opt;
opt.smoothing_scale  = 2.0 / Boxsize;     // Gaussian, box units
opt.lambda_threshold = 0.2;
opt.classification_storage = ClassificationStorage::Bytes;
CosmicWeb web;
TWeb(delta_real, opt, web);               // or TWeb(particles, NumPart, Nmesh, opt, web)
int type = web.get_classification({ix, iy, iz});   // 0 void, 1 sheet, 2 filament, 3 knot
\end{cppcode}

\file{TidalTensor/Hessian.h} computes Hessians of fields by FFT (all components from one retained Fourier copy) and their eigenvalues and
eigenvectors with GSL, for the tidal field and the shear used elsewhere (e.g.\ LPT, \cref{ch:lpt}).

\section{ORIGAMI}\label{sec:web-origami}

ORIGAMI \cite{Falck2012} classifies particles by the number of orthogonal axes along which their stream has crossed. Particles that started on a
lattice of $N^3$ points have crossed along a lattice row if their order along that direction is reversed; a particle is crossed along an axis if
it has crossed with any particle within a quarter of the row. Nine axes are checked ($x$, $y$, $z$ and the two diagonals in each coordinate
plane), grouped into four orthogonal triads, and the morphology is the largest number of crossed axes in one triad: 0 void, 1 wall, 2 filament,
3 halo. The lattice coordinate comes from the particle ID (the numbering of \code{create_particle_grid} and the initial conditions); with MPI the
rows are moved to the tasks as needed.

\code{ORIGAMIHalos} groups the halo particles on the Delaunay tessellation of all particles (\cref{sec:voids-dtfe}): Voronoi densities of all
particles, cores of Delaunay-connected halo particles above a density threshold, and the other halo particles attached in rounds to the halo of
their densest neighbour already in a halo (independent of order and number of tasks); groups without a core can be kept; no unbinding is done.

\section{Minkowski functionals}\label{sec:web-minkowski}

The morphology of an excursion set $\{\bm x:\delta(\bm x)\ge\nu\}$ is described by the four Minkowski functionals in 3D: volume, surface area,
integrated mean curvature and Euler characteristic \cite{SchmalzingBuchert1997}. \file{VoxelMorphology.h} computes them exactly for the union of
closed cubic voxels above the threshold by counting the occupied cubes $C$ and the unique faces $F$, edges $E$ and vertices $V$ of their cell
complex. With cell size $h$:
\begin{gather}
V_{\rm vol}=Ch^3,\qquad A=(2F-6C)h^2,\qquad\int H\,\dd A=\pi(3C-2F+E)h,\qquad\chi=V-E+F-C,
\end{gather}
and the conventional isotropic Crofton estimates of the underlying smooth field \cite{Schmalzing1999},
\begin{equation}
\left\{\frac{Ch^3}{L^3},\ \frac29\frac{(F-3C)h^2}{L^3},\ \frac29\frac{(E-2F+3C)h}{L^3},\ \frac{V-E+F-C}{L^3}\right\}.
\end{equation}
The two normalisations of area and curvature differ (the staircase area is $3/2$ of the Crofton one). Periodic or open boundaries are supported.
Each entity is counted exactly once by vertex-anchored patterns of the incident voxels, so the integer counts are exact across tasks and threads;
MPI exchanges one slice and reduces four counts.

\begin{validation}[Validation: cosmic web and morphology]
\file{tests/test_cosmicweb.cpp}: analytic tensors and thread-count consistency; \file{tests/test_cosmicweb_pencil.cpp}: slabs equal pencils.
Morphology: the counts against an explicit set-union enumeration; empty and full boxes, isolated cubes, face/edge/corner contacts, a cavity, a
periodic ring; uneven and empty slabs. For Gaussian random fields ($128^3$, smoothing of 4 cells, 4 realisations) the Crofton densities agree with
the continuum predictions \cite{SchmalzingBuchert1997} to $\sim1\%$ for $V_0$, $V_1$, 1--4\% for $V_2$ and 2--7\% for $V_3$ at $\nu=-2$--2, within
the statistical errors of 1--4\%.
\end{validation}
